% SPDX-License-Identifier: LicenseRef-FLOW-NC-Research-1.0
% Generated from manuscript.md by scripts/build_manuscript_tex.sh. Do not edit by hand.
% Options for packages loaded elsewhere
\PassOptionsToPackage{unicode}{hyperref}
\PassOptionsToPackage{hyphens}{url}
\documentclass[
]{article}
\usepackage{xcolor}
\usepackage{amsmath,amssymb}
\setcounter{secnumdepth}{-\maxdimen} % remove section numbering
\usepackage{iftex}
\ifPDFTeX
  \usepackage[T1]{fontenc}
  \usepackage[utf8]{inputenc}
  \usepackage{textcomp} % provide euro and other symbols
\else % if luatex or xetex
  \usepackage{unicode-math} % this also loads fontspec
  \defaultfontfeatures{Scale=MatchLowercase}
  \defaultfontfeatures[\rmfamily]{Ligatures=TeX,Scale=1}
\fi
\usepackage{lmodern}
\ifPDFTeX\else
  % xetex/luatex font selection
\fi
% Use upquote if available, for straight quotes in verbatim environments
\IfFileExists{upquote.sty}{\usepackage{upquote}}{}
\IfFileExists{microtype.sty}{% use microtype if available
  \usepackage[]{microtype}
  \UseMicrotypeSet[protrusion]{basicmath} % disable protrusion for tt fonts
}{}
\makeatletter
\@ifundefined{KOMAClassName}{% if non-KOMA class
  \IfFileExists{parskip.sty}{%
    \usepackage{parskip}
  }{% else
    \setlength{\parindent}{0pt}
    \setlength{\parskip}{6pt plus 2pt minus 1pt}}
}{% if KOMA class
  \KOMAoptions{parskip=half}}
\makeatother
\usepackage{longtable,booktabs,array}
\usepackage{calc} % for calculating minipage widths
% Correct order of tables after \paragraph or \subparagraph
\usepackage{etoolbox}
\makeatletter
\patchcmd\longtable{\par}{\if@noskipsec\mbox{}\fi\par}{}{}
\makeatother
% Allow footnotes in longtable head/foot
\IfFileExists{footnotehyper.sty}{\usepackage{footnotehyper}}{\usepackage{footnote}}
\makesavenoteenv{longtable}
\setlength{\emergencystretch}{3em} % prevent overfull lines
\providecommand{\tightlist}{%
  \setlength{\itemsep}{0pt}\setlength{\parskip}{0pt}}
\usepackage{bookmark}
\IfFileExists{xurl.sty}{\usepackage{xurl}}{} % add URL line breaks if available
\urlstyle{same}
\hypersetup{
  hidelinks,
  pdfcreator={LaTeX via pandoc}}

\author{}
\date{}

\begin{document}

\section{Individually optimal routing and total travel time on synthetic
grid networks: a preregistered simulation
study}\label{individually-optimal-routing-and-total-travel-time-on-synthetic-grid-networks-a-preregistered-simulation-study}

\textbf{Status.} Independent research manuscript / preprint draft, not
peer-reviewed. Version M4 (2026-11-25).

\subsection{Author and affiliation
metadata}\label{author-and-affiliation-metadata}

Mohamed Bechir Kefi (no affiliation; owner-approved metadata, freeze
record C3). An AI system (Claude, Anthropic) was used as a research and
development assistant; it is not an author (see the AI assistance
statement).

\subsection{Abstract}\label{abstract}

When every traveller picks the route that is fastest for themselves,
total travel time can exceed what a coordinated assignment would
achieve. We measured this gap in a small, fully offline simulation study
whose confirmatory protocol was frozen before any confirmatory instance
was solved. On 100 seeded 6 by 6 lattice networks with quartic (BPR)
link costs, at twice the demand at which free-flow routing first
saturates a link, the median excess of user-equilibrium total travel
time over the system optimum was 0.0232 (one-sided 95 percent bootstrap
lower bound 0.0213), rejecting the preregistered null hypothesis of an
excess of at most 1 percent. The decision held at solver tolerances of
1e-4 and 1e-8 and under an alternative demand calibration, but not with
linear-in-flow costs (median 0.0030) or at light demand. Free-flow
shortest-path routing was far worse than equilibrium routing under
congestion (paired median difference 0.201). In one operationalisation
of partial compliance, routing more travellers system-optimally lowered
the excess at every step. On one real test network (Sioux Falls), which
its maintainers do not consider realistic, the excess was 0.0145, just
above the margin. Two exact checks on the Braess network were
contradicted under their frozen 1e-6 rule, with residuals of 1e-5 to
6e-5 at the frozen solver tolerance; an exploratory rerun at a tighter
tolerance met the rule. The results concern static deterministic
assignment on synthetic networks and say nothing about any real city or
navigation product.

\subsection{Keywords}\label{keywords}

traffic assignment; user equilibrium; system optimum; price of anarchy;
Braess paradox; preregistration

\subsection{1. Introduction}\label{introduction}

Route-guidance tools and individual drivers usually seek the route that
is fastest for the traveller. A classical result of transport economics
is that a population of such choices, at equilibrium, need not minimise
the total time spent travelling by everyone, and that adding a road can
even make every traveller slower (Braess, Nagurney and Wakolbinger
2005). How large the gap is in a given family of networks is an
empirical question about that family.

This study asks one question: \textbf{can individually optimal routing
worsen collective travel time, and by how much, in a controlled family
of synthetic congestible networks?} It is a simulation study. It does
not model any commercial navigation product, makes no prediction for any
real city and says nothing about how real drivers behave.

Hypotheses were fixed before any confirmatory instance was solved
(Section 4.1):

\begin{itemize}
\tightlist
\item
  \textbf{Primary (EXP-003).} H0: the median over 6 by 6 grid instances
  of the relative excess \texttt{R(UE)} of user-equilibrium (UE) over
  system-optimal (SO) total travel time, at demand level rho = 2, is at
  most delta = 0.01. H1: it exceeds 0.01.
\item
  \textbf{Exact checks (EXP-001, EXP-002).} The engine reproduces the
  published Braess example and a hand-derived demand interval in which
  the extra link raises equilibrium travel time.
\item
  \textbf{Compliance (EXP-004).} The median excess decreases strictly as
  the share of travellers routed system-optimally rises.
\item
  \textbf{Braess links (EXP-005, optional).} Some 4 by 4 grid instances
  contain a link whose removal lowers equilibrium total travel time.
\item
  \textbf{Real network (EXP-006).} The primary direction also holds on
  the Sioux Falls test network.
\end{itemize}

\textbf{Contributions,} stated narrowly:

\begin{enumerate}
\def\labelenumi{\arabic{enumi}.}
\tightlist
\item
  a small tested engine that reproduces the Braess example and the
  published Sioux Falls equilibrium;
\item
  a preregistered, 100-instance measurement of the UE-to-SO gap and of
  the free-flow baseline gap on seeded BPR grids, with paired bootstrap
  intervals and sensitivity analyses that could have exposed the gap as
  an artefact of solver tolerance, cost curvature or demand calibration;
\item
  a compliance sweep and a single-link removal scan on the same
  instances;
\item
  a complete public record of every result, including two exact checks
  that failed their own frozen rule.
\end{enumerate}

No novelty is claimed for any theoretical statement.

\subsection{2. Related work}\label{related-work}

Only sources actually opened are used for content statements. Their
reading level is given in brackets; sources whose text was not opened
are cited for attribution only.

\begin{itemize}
\tightlist
\item
  \textbf{The paradox.} Braess's 1968 paper, read in the English
  translation of Braess, Nagurney and Wakolbinger (2005) {[}read{]},
  gives a four-node network in which, at total flow 6, the equilibrium
  with an added link has path flows 2, 2 and 2 and travel time 92, while
  eliminating the link makes the equilibrium coincide with the optimal
  flow. The Supernetworks Center page (Nagurney, n.d.) {[}read{]} gives
  the link costs \texttt{10f\_a}, \texttt{f\_b\ +\ 50},
  \texttt{f\_c\ +\ 50}, \texttt{10f\_d}, \texttt{f\_e\ +\ 10} and the
  travel time 83 without the added link. EXP-001 reproduces this
  example.
\item
  \textbf{Price of anarchy.} Roughgarden and Tardos (2002) {[}author
  preprint, abstract and introduction read{]} show that with linear
  latencies the total latency of selfish routing is at most 4/3 of the
  minimum, and that with general continuous nondecreasing latencies it
  can be arbitrarily worse. Our primary metric \texttt{R(UE)\ +\ 1} is
  this ratio, measured instead of bounded.
\item
  \textbf{Empirical gaps.} Youn, Gastner and Jeong (2008) {[}abstract
  read{]} study the price of anarchy on city road networks and report
  that closing certain streets can reduce travel time. We make no city
  claim; EXP-005 asks the analogous question on synthetic grids.
\item
  \textbf{Formulations and methods.} The user-equilibrium and
  system-optimum principles are commonly attributed to Wardrop (1952),
  the convex equilibrium formulation to Beckmann, McGuire and Winsten
  (1955), and the Frank-Wolfe method to Frank and Wolfe (1956) and, for
  traffic assignment, LeBlanc, Morlok and Pierskalla (1975) {[}all four:
  metadata verified, text not read; attribution only{]}. Partial central
  control and mixed behaviour classes, related to our compliance arm,
  are studied by Roughgarden (2004) and Harker (1988) {[}metadata
  verified, text not read{]}; we do not claim that our MIX(c) arm
  matches either formulation.
\item
  \textbf{Data.} The Sioux Falls network of the Transportation Networks
  for Research collection {[}README files read{]} is described by its
  maintainers as ``not considered as a realistic one'', useful ``for
  code debugging'', and as carrying the ``traditional'' BPR value B =
  0.15.
\end{itemize}

Relation to prior work: EXP-001 is a replication of a published example;
EXP-002 extends it to a demand interval derived by the assistant (not
taken from a source); EXP-003 to EXP-005 are measurements on a synthetic
family that, to our knowledge from a non-systematic search, has not been
reported in this form, which is not evidence that it has not been.

\subsection{3. Background and formal
model}\label{background-and-formal-model}

A directed graph \texttt{G\ =\ (V,\ E)} carries, on each link
\texttt{e}, a nondecreasing nonnegative travel-time function
\texttt{t\_e(x\_e)} of the link flow \texttt{x\_e}. Grid links use the
BPR form

\texttt{t\_e(x)\ =\ t0\_e\ (1\ +\ alpha\ (x\ /\ c\_e)\^{}beta)}, with
\texttt{alpha\ =\ 0.15}, \texttt{beta\ =\ 4},

where \texttt{t0\_e} is the free-flow time (minutes) and \texttt{c\_e}
the capacity (veh/h). The Braess network uses the linear costs above in
abstract units. Origin-destination (OD) pairs \texttt{w} have demand
\texttt{d\_w}; path flows \texttt{h\_p\ \textgreater{}=\ 0} satisfy
\texttt{sum\_\{p\ in\ P\_w\}\ h\_p\ =\ d\_w}, and
\texttt{x\_e\ =\ sum\_p\ delta\_\{ep\}\ h\_p}.

\begin{itemize}
\tightlist
\item
  \textbf{Total system travel time}
  \texttt{TSTT(x)\ =\ sum\_e\ x\_e\ t\_e(x\_e)} (veh min per hour).
\item
  \textbf{Primary metric}
  \texttt{R(arm)\ =\ TSTT(arm)\ /\ TSTT(SO)\ -\ 1}, dimensionless, on
  the same instance and demand level.
\item
  \textbf{FF (baseline).} All demand of each OD pair on its free-flow
  shortest path, ties broken to the lexicographically smallest node
  sequence.
\item
  \textbf{UE.} No traveller can lower their own time by switching path
  (minimises the Beckmann objective
  \texttt{sum\_e\ integral\_0\^{}\{x\_e\}\ t\_e(s)\ ds}).
\item
  \textbf{SO.} Minimises TSTT (the same solver on marginal costs
  \texttt{t\_e(x)\ +\ x\ t\_e\textquotesingle{}(x)}).
\item
  \textbf{MIX(c).} A fraction \texttt{c} of every OD demand follows
  \texttt{c} times the SO path flows; the rest reaches UE with that flow
  as fixed background. MIX(0) = UE and MIX(1) = SO by construction.
\item
  \textbf{Relative gap}
  \texttt{RG\ =\ (sum\_e\ x\_e\ t\_e(x\_e)\ -\ sum\_w\ d\_w\ kappa\_w)\ /\ sum\_e\ x\_e\ t\_e(x\_e)},
  where \texttt{kappa\_w} is the current shortest-path time of OD
  \texttt{w} (marginal costs for SO).
\item
  \textbf{Demand level rho.} For each grid instance, \texttt{D\_ref} is
  the total demand at which free-flow all-or-nothing loading first
  brings a link to \texttt{x\ /\ c\ =\ 1}. Total demand is
  \texttt{rho\ *\ D\_ref}, split over OD pairs by the instance's
  Dirichlet weights.
\end{itemize}

Two analytic facts frame the results. In the Braess network at demand 6
the UE path time is 83 without the extra link and 92 with it. For the
same network the assistant derived, and exact rational tests confirmed
before the freeze, that the link raises UE time exactly for
\texttt{d\ in\ (80/31,\ 80/9)}, with a regime change at
\texttt{d\ =\ 40/11} where the outer paths start to carry flow.

\subsection{4. Methods}\label{methods}

\subsubsection{4.1 Design and protocol
registration}\label{design-and-protocol-registration}

The study protocol (\texttt{research/protocol.md}, v1.0) and six
experiment protocols were frozen by the owner on 2026-11-24 at revision
\texttt{ca8b750}, with SHA-256 hashes of every protocol and confirmatory
config recorded in \texttt{research/freeze.json} and decision D-011. A
freeze guard in the runner refuses any confirmatory config whose hash
differs from the recorded one. The primary estimand, primary cell (6 by
6, rho = 2) and margin delta = 0.01 were written in the first draft of
the protocol (commit \texttt{91ecb92}), before any exploratory run.
Amendments made before the freeze (solver of record, sample size,
dropped demand level, real-data calibration) are listed in protocol
section 15.

\textbf{Disclosure (C4).} Before the freeze the assistant saw
exploratory values from the analogue of the primary cell on disjoint
exploratory seeds (median \texttt{R(UE)} 0.02349, n = 20; listed in
\texttt{research/freeze.json}). These informed the sample-size choice
only; delta, the primary cell and the decision rule were not changed
after they were seen.

Confirmatory work is EXP-001 to EXP-006 as frozen. Everything labelled
\textbf{exploratory} in this paper was defined after confirmatory
results existed and changes no frozen verdict.

\subsubsection{4.2 Networks and demand}\label{networks-and-demand}

\begin{itemize}
\tightlist
\item
  \textbf{Braess network} (EXP-001, EXP-002): four nodes, the five
  published links, demand 6 (EXP-001) or 75 demand values from 0.25 to
  12.0 (EXP-002).
\item
  \textbf{Grid generator} (grid-v1): \texttt{n} by \texttt{n} lattice
  with bidirectional links between 4-neighbours; per undirected pair
  \texttt{t0\ \textasciitilde{}\ U{[}1,\ 3{]}} minutes and
  \texttt{c\ \textasciitilde{}\ U{[}600,\ 1800{]}} veh/h; 12 OD pairs
  drawn without replacement among node pairs at Manhattan distance at
  least \texttt{n\ -\ 1}, weights Dirichlet(1, \ldots, 1). Confirmatory
  seeds 1000-1099 (4 by 4) and 2000-2099 (6 by 6).
\item
  \textbf{Sioux Falls} (EXP-006): the TNTP files at a pinned commit,
  checked against SHA-256 values before every use, used locally and not
  redistributed. All 76 links carry \texttt{B\ =\ 0.15} and power 4,
  which set the grid's BPR parameters (rule fixed before the fit; the
  fitted values equal the conventional ones). The published OD table
  corresponds to rho\_SF = 5.788.
\end{itemize}

\subsubsection{4.3 Routing policies and
solvers}\label{routing-policies-and-solvers}

UE, SO and the MIX follower are solved by a path-based
gradient-projection method implemented for this study (unit step; one
iteration is one Gauss-Seidel sweep over OD pairs), stopping at
\texttt{RG\ \textless{}=\ 1e-6} or 20,000 sweeps. Frank-Wolfe with exact
line search was the solver of the draft protocol; in pre-freeze
exploratory runs (journal J-0002 to J-0008) it failed to reach
\texttt{RG\ \textless{}=\ 1e-6} within 20,000 iterations on most
congested grid solves and on corner Braess solutions, so the owner
adopted gradient projection before the freeze (amendment A1) and
Frank-Wolfe was kept only as a differential reference in sensitivity S1.
A MIX(c) record counts as converged only if both its SO reference and
its follower UE converged.

\subsubsection{4.4 Metrics}\label{metrics}

Primary: \texttt{R}. Secondary: TSTT, mean trip time
(\texttt{TSTT\ /\ sum\_w\ d\_w}), demand-weighted p95 path time,
flow-weighted \texttt{x/c} and share of links over capacity, the
demand-weighted Gini coefficient of the stretch
\texttt{time\_p\ /\ FFmin\_w}, and convergence (sweeps, final RG,
wall-clock). The Braess effect is
\texttt{Delta(d)\ =\ TSTT\_UE(with\ e)\ -\ TSTT\_UE(without\ e)}.

\subsubsection{4.5 Statistical analysis}\label{statistical-analysis}

The independent unit is the instance; all arms and levels of an instance
are solved on the same network and OD sample (matched comparison).
Intervals are percentile bootstrap intervals over instances (10,000
resamples, seed 20261124), resampling instance ids for paired contrasts.
EXP-003 rejects H0 if the one-sided 95 percent lower bound of the median
exceeds 0.01. EXP-004 supports H1 if every consecutive paired median
difference has a one-sided 95 percent lower bound above 0 at capacity
scale 1.0 (no multiplicity correction; every step must pass). EXP-005
reports a Wilson 95 percent interval for the incidence. EXP-001, EXP-002
and EXP-006 are deterministic and report exact values against
tolerances, without intervals. Instances with a nonconverged UE or SO
solve are excluded and counted; if more than 10 percent of the primary
cell were excluded the primary analysis would be reported as failed.
Interpretation rules that the protocol left implicit (absolute reading
of the 1e-6 tolerance; boundary estimation by line fits; one-sided
EXP-004 bounds) were fixed in decision D-012 before any confirmatory
output.

\subsubsection{4.6 Implementation and
reproducibility}\label{implementation-and-reproducibility}

Python 3.12.1 and numpy 2.5.4 on one Apple M1 Pro machine, single
process per experiment. Every run records its config hash, git revision,
environment and per-record status. Confirmatory runs used clean revision
\texttt{8dfc66c}; the analysis code (\texttt{src/flow/confirmatory.py})
was committed before any confirmatory output (\texttt{5108a16}). The
test suite (shortest paths against brute-force enumeration, flow
conservation, nonnegative costs, disconnected OD pairs, convergence
tolerances, MIX identities, exact Braess values) passed before the runs.
Reproduction commands are in each experiment's \texttt{reproduction.md}.

\subsection{5. Results}\label{results}

Every number in this section is read from the committed
\texttt{analysis.json} of the experiment named, or from the exploratory
output named; the mapping is in
\texttt{research/publication/artifact\_manifest.json}. All six
confirmatory matrices completed: 10,626 records, of which 155 were
nonconverged (all in the Frank-Wolfe reference cell of EXP-003) and none
failed.

\subsubsection{5.1 Expected versus
observed}\label{expected-versus-observed}

\textbf{Table 1.} Predictions, observations and verdicts under the
frozen rules.

\begin{longtable}[]{@{}
  >{\raggedright\arraybackslash}p{(\linewidth - 6\tabcolsep) * \real{0.2500}}
  >{\raggedright\arraybackslash}p{(\linewidth - 6\tabcolsep) * \real{0.2500}}
  >{\raggedright\arraybackslash}p{(\linewidth - 6\tabcolsep) * \real{0.2500}}
  >{\raggedright\arraybackslash}p{(\linewidth - 6\tabcolsep) * \real{0.2500}}@{}}
\toprule\noalign{}
\begin{minipage}[b]{\linewidth}\raggedright
EXP
\end{minipage} & \begin{minipage}[b]{\linewidth}\raggedright
Prediction (frozen)
\end{minipage} & \begin{minipage}[b]{\linewidth}\raggedright
Observation
\end{minipage} & \begin{minipage}[b]{\linewidth}\raggedright
Verdict (frozen rule)
\end{minipage} \\
\midrule\noalign{}
\endhead
\bottomrule\noalign{}
\endlastfoot
EXP-001 & All Braess values within 1e-6 (absolute) of the analytic ones
& 5 of 6 arm-network combinations exact (difference at most 4.4e-16); UE
with link e: TSTT 551.99993809 against 552 (difference 6.19e-5) &
\textbf{Contradicted} \\
EXP-002 & Paradox exactly on (80/31, 80/9); boundaries within 1e-6 &
Sign of \texttt{Delta} agrees at 75 of 75 points; lower boundary within
1.3e-15; upper and middle boundaries off by 1.11e-5 and 8.76e-6 &
\textbf{Contradicted} \\
EXP-003 & Median \texttt{R(UE)} (6 by 6, rho = 2) \textgreater{} 0.01 by
its one-sided lower bound & Median 0.02317; one-sided lower bound
0.02134; n = 100, none excluded & \textbf{H0 rejected} \\
EXP-004 & Median \texttt{R(MIX(c))} strictly decreasing in c & Medians
0.02317, 0.01824, 0.01254, 0.00496, 0; all four lower bounds above 0 &
\textbf{H1 supported} \\
EXP-005 & Incidence of Braess links \textgreater{} 0 & 77 of 100
instances (Wilson 95\% 0.678 to 0.842) & \textbf{H1 supported} \\
EXP-006 & Engine gate \textless= 1e-3; \texttt{R(UE)} at rho = 2
\textgreater{} 0.01 & Gate 1.28e-4; \texttt{R(UE)} 0.01453 &
\textbf{Gate passed; direction replicates} \\
\end{longtable}

\subsubsection{5.2 Exact checks on the Braess network (EXP-001,
EXP-002)}\label{exact-checks-on-the-braess-network-exp-001-exp-002}

\textbf{EXP-001.} Free-flow and system-optimal values, and the UE
without the extra link, matched the analytic values exactly (largest
difference 4.4e-16). The UE with the link stopped after 6 sweeps at RG
7.29e-8 with TSTT 551.99993809, mean trip time 91.99998968 and zigzag
path flow 1.99999845 against 552, 92 and 2. The 6.19e-5 TSTT difference
exceeds the frozen absolute tolerance of 1e-6, so EXP-001 is
\textbf{contradicted} under its frozen rule. The published effect (83
without the link, 92 with it) is reproduced to the solver's tolerance.

\textbf{EXP-002.} Over 75 demand values, adding the link raised UE
travel time at exactly the predicted points (main sweep: 2.75 to 8.75)
and at none outside. The lower paradox boundary was recovered as
2.580645161290321 (80/31 within 1.3e-15). The two boundaries estimated
from the middle regime, where UE uses three paths, were 8.888899966
(80/9, difference 1.11e-5) and 3.636354878 (40/11, difference 8.76e-6).
EXP-002 is \textbf{contradicted} under its frozen rule. The largest
per-point UE time error was 1.06e-3, at d = 3.64 just above the regime
change (Figure
\texttt{research/experiments/EXP-002/figures/exp002\_sweep.png}).

\textbf{Protocol and config mismatch (disclosed).} EXP-001 protocol item
5 states a relative gap of 1e-10 ``for this run'' and item 17 a 1e-6
versus 1e-10 sensitivity, but the frozen config runs at 1e-6, and
EXP-002's frozen config also uses 1e-6 while its contradiction threshold
is 1e-6. The confirmatory runs used the frozen configs. The owner
decided on 2026-11-25 that the frozen verdicts stand and that a 1e-10
rerun be reported beside them as an exploratory check.

\textbf{Exploratory check at RG \textless= 1e-10 (J-0022).} The frozen
configs were rerun with only the tolerance changed and analysed with the
frozen analysis code unchanged:

\begin{longtable}[]{@{}
  >{\raggedright\arraybackslash}p{(\linewidth - 6\tabcolsep) * \real{0.2500}}
  >{\raggedright\arraybackslash}p{(\linewidth - 6\tabcolsep) * \real{0.2500}}
  >{\raggedright\arraybackslash}p{(\linewidth - 6\tabcolsep) * \real{0.2500}}
  >{\raggedright\arraybackslash}p{(\linewidth - 6\tabcolsep) * \real{0.2500}}@{}}
\toprule\noalign{}
\begin{minipage}[b]{\linewidth}\raggedright
Run (exploratory)
\end{minipage} & \begin{minipage}[b]{\linewidth}\raggedright
Records
\end{minipage} & \begin{minipage}[b]{\linewidth}\raggedright
Largest final RG
\end{minipage} & \begin{minipage}[b]{\linewidth}\raggedright
Result of the frozen analysis
\end{minipage} \\
\midrule\noalign{}
\endhead
\bottomrule\noalign{}
\endlastfoot
X-EXP-001-TOL1E-10 & 6 ok & 1.77e-11 & ``reproduced''; largest
difference 7.53e-9 (TSTT 552.0000000075) \\
X-EXP-002-TOL1E-10 & 450 ok & 9.59e-11 & ``derivation confirmed'';
boundaries within 1.3e-15 (80/31), 7.0e-10 (40/11), 1.5e-9 (80/9); sign
agrees at 75 of 75 points \\
\end{longtable}

The largest per-point UE time error in the exploratory EXP-002 run was
still 3.37e-6, at d = 3.64; per-point time agreement is not a criterion
of the frozen rule. This check is consistent with the reading that the
confirmatory contradictions come from the frozen solver tolerance rather
than from the engine or the derivation. It was defined after the
confirmatory results were seen and does not change either frozen
verdict.

\subsubsection{5.3 Primary result: equilibrium excess on grids
(EXP-003)}\label{primary-result-equilibrium-excess-on-grids-exp-003}

In the primary cell (6 by 6, rho = 2, n = 100, none excluded) the median
\texttt{R(UE)} was \textbf{0.02317} (mean 0.02575; quartiles 0.01732 to
0.03160; range 0.00386 to 0.06667). The one-sided 95 percent lower bound
was 0.02134 and the two-sided 95 percent interval {[}0.02075,
0.02559{]}. \textbf{H0 is rejected.} The upper bound is above 0.01, so
the loss is not ``practically negligible'' in the protocol's sense.

\textbf{Table 2.} Median \texttt{R} and 95 percent bootstrap interval by
grid size and demand level (n = 100 per cell; every arm converged in
every cell). Source: \texttt{research/experiments/EXP-003/analysis.json}
(\texttt{per\_cell}).

\begin{longtable}[]{@{}
  >{\raggedright\arraybackslash}p{(\linewidth - 8\tabcolsep) * \real{0.2000}}
  >{\raggedright\arraybackslash}p{(\linewidth - 8\tabcolsep) * \real{0.2000}}
  >{\raggedright\arraybackslash}p{(\linewidth - 8\tabcolsep) * \real{0.2000}}
  >{\raggedright\arraybackslash}p{(\linewidth - 8\tabcolsep) * \real{0.2000}}
  >{\raggedright\arraybackslash}p{(\linewidth - 8\tabcolsep) * \real{0.2000}}@{}}
\toprule\noalign{}
\begin{minipage}[b]{\linewidth}\raggedright
Size
\end{minipage} & \begin{minipage}[b]{\linewidth}\raggedright
rho
\end{minipage} & \begin{minipage}[b]{\linewidth}\raggedright
R(UE)
\end{minipage} & \begin{minipage}[b]{\linewidth}\raggedright
R(FF)
\end{minipage} & \begin{minipage}[b]{\linewidth}\raggedright
Paired R(FF) - R(UE)
\end{minipage} \\
\midrule\noalign{}
\endhead
\bottomrule\noalign{}
\endlastfoot
4x4 & 1 & 0.00409 {[}0.00253, 0.00553{]} & 0.0081 {[}0.0063, 0.0098{]} &
0.0008 {[}0.0000, 0.0025{]} \\
4x4 & 2 & 0.03479 {[}0.02785, 0.03803{]} & 0.3324 {[}0.3054, 0.3436{]} &
0.2951 {[}0.2654, 0.3210{]} \\
4x4 & 3 & 0.04538 {[}0.04082, 0.05448{]} & 1.6884 {[}1.5706, 1.8056{]} &
1.6332 {[}1.5038, 1.7480{]} \\
6x6 & 1 & 0.00305 {[}0.00223, 0.00361{]} & 0.0053 {[}0.0044, 0.0066{]} &
0.0005 {[}0.0000, 0.0015{]} \\
6x6 & 2 & 0.02317 {[}0.02075, 0.02559{]} & 0.2275 {[}0.2078, 0.2498{]} &
0.2007 {[}0.1782, 0.2202{]} \\
6x6 & 3 & 0.03744 {[}0.03295, 0.04152{]} & 1.2372 {[}1.1569, 1.3924{]} &
1.2101 {[}1.1188, 1.3400{]} \\
\end{longtable}

The equilibrium excess grows with demand at both sizes and is below
delta at rho = 1. The free-flow baseline is close to UE at rho = 1 and
far worse once demand is congested. Equilibrium routing raised the Gini
coefficient of the stretch slightly over SO in every cell (6 by 6, rho =
2: paired median 0.00408, interval {[}0.00279, 0.00488{]}); p95
trip-time differences were small and of both signs (supplementary Table
S3). Figure:
\texttt{research/experiments/EXP-003/figures/exp003\_R\_by\_cell.png}.

\textbf{Table 3.} Sensitivity analyses on the primary cell (6 by 6, rho
= 2). Source: \texttt{analysis.json} (\texttt{sensitivity}).

\begin{longtable}[]{@{}
  >{\raggedright\arraybackslash}p{(\linewidth - 10\tabcolsep) * \real{0.1667}}
  >{\raggedright\arraybackslash}p{(\linewidth - 10\tabcolsep) * \real{0.1667}}
  >{\raggedright\arraybackslash}p{(\linewidth - 10\tabcolsep) * \real{0.1667}}
  >{\raggedright\arraybackslash}p{(\linewidth - 10\tabcolsep) * \real{0.1667}}
  >{\raggedright\arraybackslash}p{(\linewidth - 10\tabcolsep) * \real{0.1667}}
  >{\raggedright\arraybackslash}p{(\linewidth - 10\tabcolsep) * \real{0.1667}}@{}}
\toprule\noalign{}
\begin{minipage}[b]{\linewidth}\raggedright
Cell
\end{minipage} & \begin{minipage}[b]{\linewidth}\raggedright
n used (excluded)
\end{minipage} & \begin{minipage}[b]{\linewidth}\raggedright
Median R(UE)
\end{minipage} & \begin{minipage}[b]{\linewidth}\raggedright
One-sided LB
\end{minipage} & \begin{minipage}[b]{\linewidth}\raggedright
Two-sided 95\% CI
\end{minipage} & \begin{minipage}[b]{\linewidth}\raggedright
Same decision as primary
\end{minipage} \\
\midrule\noalign{}
\endhead
\bottomrule\noalign{}
\endlastfoot
Primary (RG 1e-6) & 100 (0) & 0.02317 & 0.02134 & {[}0.02075, 0.02559{]}
& n/a \\
S1, RG 1e-4 & 100 (0) & 0.02305 & 0.02131 & {[}0.02057, 0.02552{]} &
yes \\
S1, RG 1e-8 & 100 (0) & 0.02317 & 0.02134 & {[}0.02075, 0.02559{]} &
yes \\
S2, beta = 1 & 100 (0) & 0.00304 & 0.00274 & {[}0.00264, 0.00360{]} &
\textbf{no} \\
S3, absolute calibration & 100 (0) & 0.02136 & 0.01847 & {[}0.01781,
0.02369{]} & yes \\
S1, Frank-Wolfe reference (no decision) & 7 (93) & 0.01407 & 0.00981 &
{[}0.00981, 0.02904{]} & n/a \\
\end{longtable}

Against the primary values, the paired median change was -9.9e-6 at RG
1e-4 and 5.3e-12 at RG 1e-8. The Frank-Wolfe reference converged on only
7 of 100 instances (SO) and provides no independent check. Figure:
\texttt{research/experiments/EXP-003/figures/exp003\_sensitivity.png}.

\subsubsection{5.4 Compliance and capacity
(EXP-004)}\label{compliance-and-capacity-exp-004}

\textbf{Table 4.} Median \texttt{R(MIX(c))} with 95 percent bootstrap
intervals on the 100 EXP-003 6 by 6 instances at rho = 2. Capacity scale
1.0 is confirmatory; 0.8 and 1.2 are descriptive.

\begin{longtable}[]{@{}
  >{\raggedright\arraybackslash}p{(\linewidth - 6\tabcolsep) * \real{0.2500}}
  >{\raggedright\arraybackslash}p{(\linewidth - 6\tabcolsep) * \real{0.2500}}
  >{\raggedright\arraybackslash}p{(\linewidth - 6\tabcolsep) * \real{0.2500}}
  >{\raggedright\arraybackslash}p{(\linewidth - 6\tabcolsep) * \real{0.2500}}@{}}
\toprule\noalign{}
\begin{minipage}[b]{\linewidth}\raggedright
Arm
\end{minipage} & \begin{minipage}[b]{\linewidth}\raggedright
Capacity 0.8
\end{minipage} & \begin{minipage}[b]{\linewidth}\raggedright
Capacity 1.0
\end{minipage} & \begin{minipage}[b]{\linewidth}\raggedright
Capacity 1.2
\end{minipage} \\
\midrule\noalign{}
\endhead
\bottomrule\noalign{}
\endlastfoot
FF & 0.59239 {[}0.55147, 0.65693{]} & 0.22745 {[}0.20778, 0.24982{]} &
0.09698 {[}0.08765, 0.10728{]} \\
MIX(0) = UE & 0.03058 {[}0.02698, 0.03411{]} & 0.02317 {[}0.02075,
0.02559{]} & 0.01651 {[}0.01426, 0.01935{]} \\
MIX(0.25) & 0.02542 {[}0.02237, 0.02746{]} & 0.01824 {[}0.01605,
0.02124{]} & 0.01372 {[}0.01184, 0.01556{]} \\
MIX(0.5) & 0.01758 {[}0.01643, 0.01963{]} & 0.01254 {[}0.01116,
0.01400{]} & 0.00907 {[}0.00715, 0.01043{]} \\
MIX(0.75) & 0.00868 {[}0.00809, 0.00950{]} & 0.00496 {[}0.00432,
0.00600{]} & 0.00294 {[}0.00251, 0.00338{]} \\
MIX(1) = SO & 0 & 0 & 0 \\
\end{longtable}

At capacity 1.0 the four consecutive paired median differences were
0.00415, 0.00539, 0.00725 and 0.00496, with one-sided lower bounds
0.00343, 0.00442, 0.00652 and 0.00445, all above 0: \textbf{H1 is
supported}. Every step also passed at scales 0.8 and 1.2 (descriptive).
The identity checks held: MIX(0) reproduced the EXP-003 UE TSTT on all
100 instances (largest relative difference 0.0) and MIX(1) gave
\texttt{R\ =\ 0}. Compliant travellers' mean trip time exceeded that of
the others by a median factor of 1.0051 at c = 0.25 and 1.0132 at c =
0.75 (capacity 1.0). Figure:
\texttt{research/experiments/EXP-004/figures/exp004\_compliance.png}.

\subsubsection{5.5 Braess links in 4 by 4 grids (EXP-005,
optional)}\label{braess-links-in-4-by-4-grids-exp-005-optional}

On the 100 EXP-003 4 by 4 instances at rho = 2, each of the 48 directed
links was removed in turn (5,154 solves, no disconnecting removal).
\textbf{77 of 100 instances} contained at least one link whose removal
lowered TSTT\_UE by more than 1e-4 relative (Wilson 95 percent {[}0.678,
0.842{]}); all 177 such links were confirmed at RG 1e-8. At the 1e-3
threshold the incidence was 67 of 100 ({[}0.573, 0.754{]}). The largest
relative decrease was 0.0354 and the smallest flagged one 1.16e-4.
\textbf{H1 is supported.} An exploratory materiality check (J-0023)
found that 17 of 100 instances have a removal lowering TSTT\_UE by more
than 1e-2, and that the median effect over the 177 flagged links was
0.0030. Figure:
\texttt{research/experiments/EXP-005/figures/exp005\_removals.png}.

\subsubsection{5.6 One real test network
(EXP-006)}\label{one-real-test-network-exp-006}

The engine's UE link flows at the published Sioux Falls demand matched
the TNTP best-known solution within 1.28e-4 of the largest reference
flow (gate threshold 1e-3): the gate passed.

\textbf{Table 5.} Sioux Falls, deterministic, one network (no interval).
Source: \texttt{research/experiments/EXP-006/analysis.json}.

\begin{longtable}[]{@{}llll@{}}
\toprule\noalign{}
Demand level & R(UE) & R(FF) & Gini of stretch UE / SO / FF \\
\midrule\noalign{}
\endhead
\bottomrule\noalign{}
\endlastfoot
rho = 1 & 0.00463 & 0.00614 & 0.0144 / 0.0094 / 0.0153 \\
rho = 2 & \textbf{0.01453} & 0.21769 & 0.0424 / 0.0305 / 0.1866 \\
rho = 3 & 0.03012 & 1.12283 & 0.0786 / 0.0557 / 0.4639 \\
published (rho 5.788) & 0.03973 & 8.36129 & 0.1950 / 0.1774 / 0.7022 \\
\end{longtable}

At RG 1e-8, \texttt{R(UE)} was 0.01452 at rho = 2 and 0.03975 at the
published level. The direction replicates at rho = 2 by a margin of
0.0045 above delta. Figure:
\texttt{research/experiments/EXP-006/figures/exp006\_R.png}.

\subsubsection{5.7 Runtime}\label{runtime}

EXP-003, which includes the Frank-Wolfe reference, took 2,222.8 s on one
machine; every other confirmatory experiment took under one minute
(EXP-005: 48.6 s). Mean gradient-projection time per solve in the
primary cell was 0.022 s (UE) and 0.049 s (SO).

\subsection{6. Discussion}\label{discussion}

\subsubsection{6.1 What the evidence
supports}\label{what-the-evidence-supports}

In this synthetic family, with quartic BPR costs and demand at twice the
free-flow saturation level, equilibrium routing costs a median of about
2.3 percent more total travel time than the system optimum. The excess
is not carried by a few instances: in the exploratory check J-0023, 90
of the 100 primary-cell instances individually exceeded delta. It grows
with demand (Table 2) and, in the same exploratory and purely
correlational check, tracks congestion within the primary cell (Spearman
correlation 0.653 between \texttt{R(UE)} and the UE flow-weighted
volume-to-capacity ratio, n = 100). This is consistent with, but does
not establish, the mechanism named in the protocol: equilibrium
travellers ignore the marginal delay they impose on others, and that
delay matters only where links are loaded into the steep part of the
cost curve.

The free-flow baseline matters more than the equilibrium inefficiency.
At rho = 2 ignoring congestion altogether costs about nine times the
equilibrium excess (paired median \texttt{R(FF)\ -\ R(UE)} 0.201 on 6 by
6 grids). At rho = 1 all three policies are within 1 percent of each
other. A reader who equates ``individually optimal'' with ``shortest
path on an empty network'' would therefore overstate what
congestion-aware individual routing loses.

\subsubsection{6.2 Competing explanations}\label{competing-explanations}

\textbf{Table 6.} Competing explanations named in the protocols and what
the results say about them.

\begin{longtable}[]{@{}
  >{\raggedright\arraybackslash}p{(\linewidth - 6\tabcolsep) * \real{0.2500}}
  >{\raggedright\arraybackslash}p{(\linewidth - 6\tabcolsep) * \real{0.2500}}
  >{\raggedright\arraybackslash}p{(\linewidth - 6\tabcolsep) * \real{0.2500}}
  >{\raggedright\arraybackslash}p{(\linewidth - 6\tabcolsep) * \real{0.2500}}@{}}
\toprule\noalign{}
\begin{minipage}[b]{\linewidth}\raggedright
Conclusion
\end{minipage} & \begin{minipage}[b]{\linewidth}\raggedright
Competing explanation (protocol)
\end{minipage} & \begin{minipage}[b]{\linewidth}\raggedright
Addressed by
\end{minipage} & \begin{minipage}[b]{\linewidth}\raggedright
Status
\end{minipage} \\
\midrule\noalign{}
\endhead
\bottomrule\noalign{}
\endlastfoot
EXP-003 primary & Solver nonconvergence inflating or shrinking the gap &
S1: decision unchanged at RG 1e-4 and 1e-8; paired shifts below 1e-5 &
Addressed for gradient projection. No independent solver check: the
Frank-Wolfe reference failed to converge on 93 of 100 instances. \\
EXP-003 primary & Dependence on cost curvature & S2 (beta = 1): median
0.00304, below delta; 0 of 100 instances above delta (J-0023) &
\textbf{Confirmed as a scope limit.} The result holds for beta = 4
only. \\
EXP-003 primary & Artefact of the demand calibration & S3 (absolute
calibration): median 0.02136, same decision; 86 of 100 instances above
delta (J-0023) & Addressed for the two calibrations tested. The result
does depend on the demand level: below delta at rho = 1. \\
EXP-003 primary & OD sampling & 100 independent instances; 90 of 100
above delta individually (J-0023); 4 by 4 grids at rho = 2, with
different draws, also exceed delta (median 0.03479) & Partly addressed.
The OD count (12) and distance rule were never varied. \\
EXP-001, EXP-002 & Implementation bug, nonconvergence, tie-breaking,
tolerance & FF tie rule exact; exact rational tests of the derivation;
exploratory 1e-10 rerun meets the rule (J-0022) & The exploratory
evidence favours the frozen solver tolerance. A small implementation
error cannot be excluded by the confirmatory run alone. \\
EXP-004 & Re-equilibration of non-compliant users cancelling the
benefit; nonmonotone effects; follower nonconvergence & Every step
positive at every capacity scale; all MIX records converged & Addressed
within this MIX definition. The link-flow scaling sensitivity (protocol
item 17) was not run (owner decision, 2026-11-25). \\
EXP-005 & Apparent reductions within solver tolerance & All 177 flagged
links confirmed at RG 1e-8 & Addressed. Materiality is a separate
question (Section 5.5). \\
EXP-006 & Ingestion error; nonconvergence; demand scaling; single
network & Engine gate 1.28e-4; S1 at 1e-8 unchanged; R reported at four
levels & Ingestion and convergence addressed. One network supports no
population claim. \\
\end{longtable}

\subsubsection{6.3 Contradictory and weak
results}\label{contradictory-and-weak-results}

The two exact checks failed their own frozen rule. The reason is visible
in the records: the frozen configs stop the solver at a relative gap of
1e-6 while the contradiction threshold is an absolute 1e-6 on quantities
of order 500, which a solver stopping at that gap need not meet. This is
a design error in the frozen protocol and config, not a property of the
network; it is reported as a contradiction because that is what the
frozen rule gives. The exploratory 1e-10 rerun shows that the engine
meets the rule when the solver is run to the tolerance the EXP-001
protocol text intended.

The real-network replication is weak: 0.01453 against a margin of 0.01,
on one network that its maintainers do not consider realistic, at a
demand level defined by our own calibration. The compliance result
(EXP-004) is close to what the MIX construction implies: MIX(0) is UE
and MIX(1) is SO, so an overall decline is expected; the informative
part is that every intermediate step declined despite re-equilibration.
Most Braess links found by EXP-005 are small: the median effect is
0.0030 of TSTT\_UE, below the study's own practical margin.

\subsubsection{6.4 Practical significance}\label{practical-significance}

The excess measured here is modest: about 2 percent at the primary
level, a median below 5 percent in every grid cell, and below 1 percent
with linear costs or light demand. With linear costs (S2) it is far
below the worst-case ratio of 4/3 that Roughgarden and Tardos (2002)
prove for linear latencies. The margin delta = 0.01 is a convention
fixed by the study before any run, not an externally grounded threshold
of practical importance. The study supports the statement that
equilibrium routing can worsen collective travel time by more than this
preregistered 1 percent margin in congested synthetic networks with
conventional quartic costs; it does not support any statement about the
size of the effect on real roads.

\subsection{7. Limitations and threats to
validity}\label{limitations-and-threats-to-validity}

\begin{itemize}
\tightlist
\item
  \textbf{Synthetic networks.} The primary hypothesis rests on seeded
  lattices with 12 OD pairs. Sioux Falls is one small test network. No
  transfer to other topologies, sizes or real cities is implied.
\item
  \textbf{Static deterministic assignment.} No time dynamics, queues,
  departure-time choice, signals or stochastic route choice.
\item
  \textbf{Behavioural model.} ``Individually optimal'' means free-flow
  shortest paths or Wardrop equilibrium with perfect information.
  Synthetic agents are not people; nothing is inferred about real
  drivers or any navigation product.
\item
  \textbf{Cost curvature.} The effect falls below delta at beta = 1
  (S2). The BPR parameters were ``fitted'' from a file that itself
  assumes the traditional values, so the calibration confirms convention
  rather than measuring roads.
\item
  \textbf{Demand definition.} rho is relative to free-flow saturation;
  at rho = 1 the effect is below delta.
\item
  \textbf{Solver.} All arms use one solver implementation; the
  Frank-Wolfe reference did not converge on most instances, so there is
  no independent solver check on grids. The engine is externally
  validated on the Braess network and on Sioux Falls (UE only; no
  published SO was used).
\item
  \textbf{Exact checks.} EXP-001 and EXP-002 are contradicted under
  their frozen rule; the 1e-10 evidence is exploratory and post hoc.
\item
  \textbf{Pre-freeze exposure.} Exploratory values of the primary-cell
  analogue were seen before the freeze (C4).
\item
  \textbf{Shared instances.} EXP-004 and EXP-005 reuse EXP-003
  instances; they are not independent replications.
\item
  \textbf{Untested sensitivities.} EXP-004 link-flow scaling (item 17)
  was not run, by owner decision.
\item
  \textbf{Provenance flags.} Three confirmatory manifests record
  \texttt{dirty:\ true} because untracked analysis outputs existed
  during the runs; no tracked file differed (D-014). One test executed
  the EXP-005 config once before the official run; its output was not
  inspected (D-013).
\item
  \textbf{Statistical.} Secondary contrasts are descriptive and not
  corrected for multiplicity. Nonrejection at rho = 1 is not evidence of
  equivalence.
\end{itemize}

\subsection{8. Ethics and responsible
research}\label{ethics-and-responsible-research}

The study uses synthetic data and one public research network (TNTP
Sioux Falls, ``for academic research purposes only''), used locally and
not redistributed; its source is cited. No personal data, human
participants or behavioural data were involved, so no ethics review was
required and none is claimed. The results must not be presented as
evidence about real drivers, real cities or any commercial routing
product.

\subsection{9. Conclusion and future
work}\label{conclusion-and-future-work}

On congested synthetic grids with quartic BPR costs, individually
optimal (equilibrium) routing increased total travel time over the
system optimum by a median of 2.3 percent, above the preregistered 1
percent margin, robustly to solver tolerance and demand calibration but
not to linear costs, and not at light demand (rho = 1), where all
policies were within 1 percent of the optimum in median. Free-flow
routing was far worse; partial system-optimal compliance reduced the
excess at every step; small Braess links were common; one real test
network showed the same direction at a small margin. Two exact checks
failed their frozen rule because of a protocol and config tolerance
mismatch.

Feasible extensions, each needing a new registered protocol: an
independent second solver that converges on grids; varied OD counts and
topologies; the link-flow compliance sensitivity; further TNTP networks;
and a registered rerun of the exact checks at the intended tolerance.

\subsection{Data and code
availability}\label{data-and-code-availability}

All code, configs, frozen protocols, analyses, figures and confirmatory
per-record tables are in this repository (paths given in each experiment
report). The Sioux Falls files are not redistributed;
\texttt{scripts/download\_tntp.py} fetches them from the pinned commit
and checks their SHA-256 values; the EXP-006 per-link results are kept
local for the same reason. Repository URL: {[}to be added on
owner-approved release{]}.

\subsection{Funding, competing interests and AI
assistance}\label{funding-competing-interests-and-ai-assistance}

No funding. No competing interests. An AI system (Claude, Anthropic)
acted as research and development assistant: it drafted protocols,
implemented the engine and tests, ran the experiments and drafted the
reports and this manuscript. The owner approved the protocol freeze, the
amendments and every decision recorded as an owner decision. Numbers
were checked by an automated trace against committed result files
(\texttt{research/publication/evidence\_audit.md}) and an internal
second review pass (\texttt{research/publication/review\_pass\_2.md}).
These are internal quality controls, not peer review.

\subsection{References}\label{references}

\begin{itemize}
\tightlist
\item
  Beckmann, M. J., McGuire, C. B., Winsten, C. B. (1955). \emph{Studies
  in the Economics of Transportation}. RAND Corporation, RM-1488.
  https://www.rand.org/pubs/research\_memoranda/RM1488.html {[}metadata
  verified; text not read{]}.
\item
  Braess, D., Nagurney, A., Wakolbinger, T. (2005). On a paradox of
  traffic planning. \emph{Transportation Science} 39(4), 446-450.
  doi:10.1287/trsc.1050.0127 {[}read{]}.
\item
  Frank, M., Wolfe, P. (1956). An algorithm for quadratic programming.
  \emph{Naval Research Logistics Quarterly} 3(1-2), 95-110.
  doi:10.1002/nav.3800030109 {[}metadata verified; text not read{]}.
\item
  Harker, P. T. (1988). Multiple equilibrium behaviors on networks.
  \emph{Transportation Science} 22(1), 39-46. doi:10.1287/trsc.22.1.39
  {[}metadata verified; text not read{]}.
\item
  LeBlanc, L. J., Morlok, E. K., Pierskalla, W. P. (1975). An efficient
  approach to solving the road network equilibrium traffic assignment
  problem. \emph{Transportation Research} 9(5), 309-318.
  doi:10.1016/0041-1647(75)90030-1 {[}metadata verified; text not
  read{]}.
\item
  Nagurney, A. (n.d.). The Braess Paradox. Virtual Center for
  Supernetworks, UMass Amherst.
  https://supernet.isenberg.umass.edu/facts/braess.html {[}read
  2026-11-24{]}.
\item
  Roughgarden, T. (2004). Stackelberg scheduling strategies. \emph{SIAM
  Journal on Computing} 33(2), 332-350. doi:10.1137/S0097539701397059
  {[}metadata verified; text not read{]}.
\item
  Roughgarden, T., Tardos, E. (2002). How bad is selfish routing?
  \emph{Journal of the ACM} 49(2), 236-259. doi:10.1145/506147.506153
  {[}author preprint of 5 December 2001, abstract and introduction
  read{]}.
\item
  Transportation Networks for Research Core Team. Transportation
  Networks for Research (TNTP), repository README and
  SiouxFalls/README.md at commit
  d1639b4ef218c17928ba573e806ddf8ba5e7ae6d.
  https://github.com/bstabler/TransportationNetworks {[}read 2026-11-24
  and 2026-11-25{]}.
\item
  Wardrop, J. G. (1952). Some theoretical aspects of road traffic
  research. \emph{Proceedings of the Institution of Civil Engineers}
  1(3), 325-362. doi:10.1680/ipeds.1952.11259 {[}metadata verified; text
  not read{]}.
\item
  Youn, H., Gastner, M. T., Jeong, H. (2008). Price of anarchy in
  transportation networks: efficiency and optimality control.
  \emph{Physical Review Letters} 101(12), 128701.
  doi:10.1103/PhysRevLett.101.128701 {[}abstract read{]}.
\end{itemize}

\subsection{Supplementary appendix}\label{supplementary-appendix}

See \texttt{research/paper/supplementary.md}: complete experimental
matrix (S1), extended tables (S2, S3), configurations (S4), robustness
and exploratory results (S5), experiment reports (S6) and deviations and
owner decisions (S7).

\end{document}
